% Generated draft: requires mathematical and source review.
\documentclass[11pt]{article}
\usepackage[margin=25mm]{geometry}
\usepackage{fontspec}
\setmainfont{Latin Modern Roman}
\usepackage{amsmath,amssymb,mathtools,unicode-math}
\setmathfont{Latin Modern Math}
\usepackage{enumitem,needspace,xurl,hyperref}
\setlength{\parindent}{0pt}
\setlength{\parskip}{5pt}
\setlength{\emergencystretch}{3em}
\newcommand{\N}{\mathbb{N}}
\newcommand{\Z}{\mathbb{Z}}
\newcommand{\Q}{\mathbb{Q}}
\newcommand{\R}{\mathbb{R}}
\newcommand{\C}{\mathbb{C}}
\newcommand{\F}{\mathbb{F}}
\newcommand{\A}{\mathbb{A}}
\newcommand{\PP}{\mathbb{P}}
\newcommand{\E}{\mathbb{E}}
\newcommand{\eps}{\varepsilon}
\newcommand{\dd}{\,\mathrm{d}}
\begin{document}
\section*{UnsolvedMath ID: 700504}
\textbf{Source problem code:} catalogue entry
\subsection{Definitions and mathematical statement}
A group \(G\) is metabelian if \(G''=[G',G']=1\), where \(G'=[G,G]\) is its commutator subgroup. A group variety is the class of all groups satisfying a specified set of group identities. A metabelian variety is a group variety all of whose members are metabelian.

For each integer \(n\geq 1\), let \(F_n\) denote the free group on \(n\) generators. An identity of a group \(G\) is an equation \(w(x_1,\ldots,x_n)=1\), with \(w\in F_n\), that holds for every tuple \((g_1,\ldots,g_n)\in G^n\). Write \(\operatorname{Id}(G)\) for the collection of all such identities, over all \(n\geq 1\), and write \(\operatorname{Var}(G)\) for the variety of all groups satisfying every identity in \(\operatorname{Id}(G)\). Here a group is called \(2\)-generated if \(G=\langle a,b\rangle\) for some \(a,b\in G\); it need not require two generators minimally.

The problem proposed by A. Yu. Olshanskii in 1982 asks:
\begin{enumerate}
\item[(a)] Classify all metabelian varieties of groups, or show that this classification is, in a suitable sense, a ``wild'' problem. The source calls this a well-known problem, but specifies neither a formal meaning of ``wild'' nor a required form of classification.
\item[(b)] Describe the identities of \(2\)-generated metabelian groups: classify the possible collections \(\operatorname{Id}(G)\) as \(G\) ranges over groups satisfying
\[
G=\langle a,b\rangle\quad\text{for some }a,b\in G,
\qquad G''=1.
\]
Equivalently, classify the distinct varieties \(\operatorname{Var}(G)\) arising from these groups.
\end{enumerate}
The formulation of (b) above interprets the source's phrase ``varieties generated by such groups'' as varieties generated by an individual such group. If arbitrary collections \(\mathcal C\) of \(2\)-generated metabelian groups are intended instead, the target would be the varieties defined by the common identities \(\bigcap_{G\in\mathcal C}\operatorname{Id}(G)\); the source does not explicitly distinguish these readings.

The supplied registry labels the problem open. It also reports an explicitly unverified AI classification claim for part (b); that report is not a verified resolution.
\subsection{Short English statement}
Classify varieties consisting of metabelian groups, or make precise and establish the proposed wildness of that classification. Also describe which identity collections occur in metabelian groups generated by two elements, equivalently which varieties those individual groups generate. The meaning of wildness and the possible collection-based reading of the second request require clarification. The supplied claim of a classification in part (b) is explicitly unverified.
\subsection{Sources}
\begin{itemize}
\item Local UnsolvedMath registry, supplied complete source record: permanent ID \(700504\), exact source problem code \texttt{catalogue entry}, title \emph{Kourovka Notebook Problem 8.54}; upstream revision \texttt{d371ee917ac9647fbb06cd6fc529deec6f3dda82}.
\item \emph{The Kourovka Notebook (Unsolved Problems in Group Theory)}, problem 8.54, proposed by A. Yu. Olshanskii, 8th issue (1982); supplied reference to the 21st issue edition (2026), \texttt{arXiv:1401.0300v48}. The registry describes the collection as edited by E. I. Khukhro and V. D. Mazurov. This attribution and reference are supplied by the record and have not been independently checked.
\item The supplied discussion attributes an explicitly unverified AI claim that the varieties in part (b) are classified to Achyuth Jayadevan, \emph{Preprint} (2026), \url{https://zenodo.org/records/22843496}. The supplied status review says this claim is marked with \(\#\) in the October 2026 version of the 21st issue. No verification of the claim is supplied.
\item Supplied status-review date: \texttt{2026-10-16}; canonical status: open; evidence category: \texttt{unrefereed\_claim}. This is supplied dated curation, not an independently verified assessment of the literature or mathematical status.
\end{itemize}
\end{document}
